% ============================================================================
% MÓDULO 3 - Visualización de datos con ggplot2
% ============================================================================
\chapter{Visualización de datos con ggplot2}
\label{cap:visualizacion}

\begin{objetivos}
\begin{itemize}
  \item Explicar por qué una gráfica revela cosas que una tabla de
        estadísticas esconde (el cuarteto de Anscombe).
  \item Construir cualquier gráfica con la \emph{gramática de gráficos} de
        \paquete{ggplot2}: datos, estéticas, geometrías, escalas, facetas,
        coordenadas y temas, capa por capa.
  \item Elegir la gráfica correcta según la pregunta de negocio:
        comparación, evolución, distribución, relación o composición.
  \item Crear barras, líneas, dispersión, histogramas, boxplots, mapas de
        calor y facetas con formato de negocio (\$, \%, miles, millones).
  \item Aplicar reglas profesionales de diseño: un color por serie,
        resaltar lo importante, sin doble eje, sin arcoíris, ejes honestos.
  \item Construir tu propio tema corporativo \codigo{tema\_empresa()} y
        exportar gráficas listas para PowerPoint, Word o correo.
  \item Armar un tablero de ventas coherente para la dirección y
        redactar conclusiones basadas en lo que los datos muestran.
\end{itemize}
\end{objetivos}

Es jueves por la tarde y la directora general te pide ``tres o cuatro
láminas'' con las ventas de 2023 para el consejo del lunes. En el
Módulo~\ref{cap:manipulacion} aprendiste a obtener las cifras con
\paquete{dplyr}: ingresos por mes, por categoría, por tienda. El problema es
que nadie en el consejo va a leer una tabla de 12 filas y 5 columnas. Lo que
sí van a hacer es mirar una gráfica durante diez segundos y quedarse con una
idea. Tu trabajo es que esa idea sea la correcta.

Visualizar datos no es ``hacer dibujos bonitos''. Es la parte del
\termino{Business Intelligence} (BI) donde el análisis se convierte en una
decisión: una gráfica bien hecha hace evidente que abril fue un mes malo, que
una tienda pequeña está vendiendo más que la grande o que una categoría se
desplomó en el segundo semestre. Una gráfica mal hecha (colores sin sentido,
ejes recortados, pasteles con diez rebanadas) confunde o, peor, engaña.

En este módulo aprenderás \paquete{ggplot2}, el paquete de gráficas más usado
en R, y, sobre todo, aprenderás a \emph{pensar} como diseñador de
visualizaciones de negocio. Cada gráfica seguirá la misma ruta: una pregunta
de negocio, el código comentado línea por línea, la figura resultante, cómo
leerla y un reto para que la hagas tú. Al final (sección~\ref{sec:m3-caso})
armarás un tablero de ventas 2023 completo para la dirección.

% ----------------------------------------------------------------------------
\section{Por qué visualizar}
\label{sec:m3-por-que}

\subsection{El cuarteto de Anscombe}

En 1973 el estadístico Francis Anscombe publicó cuatro pequeños conjuntos de
datos, cada uno con 11 pares de valores $(x, y)$. Vienen incluidos en R con
el nombre \codigo{anscombe}. Su gracia es que, si solo miras las
estadísticas de resumen, los cuatro parecen idénticos. Empecemos cargando
los paquetes que usaremos en todo el módulo.

\begin{rcode}
library(tidyverse)   # incluye ggplot2, dplyr, tidyr, forcats...
library(scales)      # formatos de ejes: dólares, comas, porcentajes
\end{rcode}

\paquete{ggplot2} ya viene dentro de \paquete{tidyverse}, así que no hace
falta cargarlo aparte. \paquete{scales} se instala junto con
\paquete{ggplot2} pero hay que cargarlo para usar sus funciones de formato
(\codigo{dollar()}, \codigo{percent()}\ldots). La tabla \codigo{anscombe}
tiene las columnas \codigo{x1}, \codigo{x2}, \codigo{x3}, \codigo{x4},
\codigo{y1}, \ldots, \codigo{y4}; la pasamos a formato largo (una fila por
punto) con \codigo{pivot\_longer()}, que conociste en el
Módulo~\ref{cap:manipulacion}.

\begin{rcode}
anscombe_largo <- anscombe %>%
  pivot_longer(cols = everything(),
               # ".value" = la 1a letra (x o y) será el nombre de columna
               names_to = c(".value", "conjunto"),
               # el patrón separa "x1" en "x" y "1"
               names_pattern = "(.)(.)") %>%
  mutate(conjunto = paste("Conjunto", conjunto))

# Las estadísticas de cada conjunto
anscombe_largo %>%
  group_by(conjunto) %>%
  summarise(media_x     = mean(x),
            media_y     = mean(y),
            desv_y      = sd(y),
            correlacion = cor(x, y),
            .groups = "drop") %>%
  mutate(across(where(is.numeric), ~ round(.x, 2)))
\end{rcode}
\begin{rsalida}
PENDIENTE
\end{rsalida}

Las cuatro filas son prácticamente iguales: media de $x$ igual a 9, media de
$y$ de 7.5, la misma desviación estándar y una correlación de 0.82 en todos
los casos. Si esto fuera un reporte con solo la tabla, concluirías que los
cuatro conjuntos ``se comportan igual''. Ahora grafiquémoslos.

\begin{rcode}
g_anscombe <- ggplot(anscombe_largo, aes(x = x, y = y)) +
  geom_point(color = "#2a78d6", size = 2) +        # los 11 puntos
  geom_smooth(method = "lm", formula = y ~ x,      # recta de regresión
              se = FALSE, color = "#eb6834", linewidth = 0.8) +
  facet_wrap(~ conjunto, ncol = 4) +               # un panel por conjunto
  labs(title = "Mismas estadísticas, cuatro historias distintas",
       subtitle = "Media, varianza, correlación y recta casi idénticas",
       caption = "Fuente: datasets::anscombe (Anscombe, 1973)") +
  theme_minimal() +
  theme(panel.grid.minor = element_blank())
g_anscombe
\end{rcode}

No te preocupes todavía por cada línea del código: en la siguiente sección
lo construiremos pieza por pieza. Fíjate en el resultado
(figura~\ref{fig:m03-anscombe}).

\begin{figure}[H]
\centering
\includegraphics[width=0.95\textwidth]{m03-anscombe}
\caption{El cuarteto de Anscombe: cuatro conjuntos con estadísticas
idénticas y formas completamente distintas.}
\label{fig:m03-anscombe}
\end{figure}

\textbf{Cómo leerla.} El conjunto 1 es una relación lineal con ruido: la
recta sí lo describe. El conjunto 2 es una curva: una recta es el modelo
equivocado. En el conjunto 3 todos los puntos caen sobre una recta perfecta
salvo uno, un \termino{valor atípico} (\emph{outlier}) que jala la recta
hacia arriba. En el conjunto 4 casi todos los puntos tienen el mismo $x$ y un
solo punto extremo fabrica toda la ``correlación''. La lección para un
analista de BI es directa: \textbf{antes de reportar un promedio o una
correlación, grafica los datos}.

\begin{negocio}{El promedio que escondía dos tiendas}
Una cadena reporta que el ticket promedio de sus tiendas es de \$950 y que
es ``estable''. Al graficar la distribución de tickets por tienda aparece que
una sucursal tiene muchísimas compras pequeñas y unas cuantas enormes
(clientes mayoristas). El promedio era el mismo, pero la operación, el
riesgo y la estrategia comercial son completamente distintos. Solo la
gráfica lo hizo evidente.
\end{negocio}

\subsection{Visualización exploratoria contra explicativa}

En tu trabajo harás dos tipos de gráficas muy diferentes:

\begin{itemize}
  \item \termino{Visualización exploratoria}: gráficas rápidas que haces
        \emph{para ti} mientras analizas. No necesitan título bonito ni
        colores corporativos; necesitan ser rápidas. Harás decenas y
        tirarás casi todas. El objetivo es \emph{descubrir}.
  \item \termino{Visualización explicativa}: la gráfica que llevas a la
        dirección. Tiene \emph{un} mensaje, un título que lo dice
        (``Belleza cayó 41\% en el segundo semestre'' en vez de ``Ventas por
        categoría''), formato de negocio en los ejes, la fuente de los datos
        y nada que distraiga. El objetivo es \emph{comunicar}.
\end{itemize}

\begin{consejo}
Un buen título explicativo es una \emph{conclusión}, no una descripción.
Compara ``Ingresos mensuales 2023'' con ``7 de 12 meses superaron la meta''.
El segundo le ahorra al director el trabajo de interpretar la gráfica. El
subtítulo puede decir qué se mide (``Ingresos mensuales, MXN'').
\end{consejo}

% ----------------------------------------------------------------------------
\section{La gramática de gráficos}
\label{sec:m3-gramatica}

\paquete{ggplot2} se basa en la \termino{gramática de gráficos}
(\emph{grammar of graphics}), una idea de Leland Wilkinson que Hadley Wickham
implementó en R: toda gráfica estadística se puede describir como la suma de
unas cuantas piezas independientes. En vez de elegir ``un tipo de
gráfica'' de un menú (como en Excel), \emph{construyes} la gráfica
declarando cada pieza.

\begin{figure}[H]
\centering
\begin{tikzpicture}[
  capa/.style={draw=#1, fill=#1!10, rounded corners=2pt,
               minimum width=4.2cm, minimum height=0.62cm,
               font=\sffamily\small\bfseries, text=#1},
  expl/.style={font=\sffamily\footnotesize, text=gristexto, anchor=west,
               align=left},
  fun/.style={font=\ttfamily\footnotesize, text=codigoinline, anchor=east}]
  \node[capa=primario]   (dat) at (0,0)    {1. Datos};
  \node[capa=secundario] (aes) at (0,0.75) {2. Estéticas (aes)};
  \node[capa=exito]      (geo) at (0,1.5)  {3. Geometrías (geom)};
  \node[capa=acento]     (esc) at (0,2.25) {4. Escalas (scale)};
  \node[capa=morado]     (fac) at (0,3.0)  {5. Facetas (facet)};
  \node[capa=peligro]    (coo) at (0,3.75) {6. Coordenadas (coord)};
  \node[capa=gristexto]  (tem) at (0,4.5)  {7. Tema (theme, labs)};
  \node[expl] at (2.4,0)    {La tabla: una fila por observación};
  \node[expl] at (2.4,0.75) {Qué columna va a $x$, $y$, color, tamaño\ldots};
  \node[expl] at (2.4,1.5)  {Cómo se dibuja: puntos, barras, líneas\ldots};
  \node[expl] at (2.4,2.25) {Cómo se traduce el dato: \$, \%, colores, log};
  \node[expl] at (2.4,3.0)  {Dividir en paneles pequeños};
  \node[expl] at (2.4,3.75) {Sistema de ejes: normal, invertido, zoom};
  \node[expl] at (2.4,4.5)  {Títulos, fuentes, cuadrícula, leyenda};
  \node[fun] at (-2.3,0)    {ggplot(datos)};
  \node[fun] at (-2.3,0.75) {aes(x, y)};
  \node[fun] at (-2.3,1.5)  {geom\_col()};
  \node[fun] at (-2.3,2.25) {scale\_y\_continuous()};
  \node[fun] at (-2.3,3.0)  {facet\_wrap()};
  \node[fun] at (-2.3,3.75) {coord\_cartesian()};
  \node[fun] at (-2.3,4.5)  {theme\_minimal()};
  \draw[-{Stealth[length=2.5mm]}, thick, acento] (8.9,0) -- (8.9,4.5)
     node[midway, right, font=\sffamily\scriptsize, text=acento,
          align=left] {se suman\\con \texttt{+}};
\end{tikzpicture}
\caption{Las siete capas de la gramática de gráficos. Solo las tres
primeras son obligatorias; las demás tienen valores por defecto.}
\label{fig:m3-capas}
\end{figure}

Las capas se combinan con el signo \codigo{+}. Una gráfica mínima necesita
solo las tres primeras: \emph{datos}, \emph{estéticas} y al menos una
\emph{geometría}. Todo lo demás (escalas, facetas, coordenadas, tema) tiene
un valor por defecto razonable que puedes cambiar.

\begin{concepto}{Estética (aes)}
Una \termino{estética} es una propiedad visual que \emph{depende de los
datos}: la posición en $x$, la posición en $y$, el color, el relleno
(\codigo{fill}), el tamaño, la forma, la transparencia (\codigo{alpha}).
Dentro de \codigo{aes()} dices ``esta columna controla esta propiedad''.
Si quieres que \emph{todas} las barras sean azules, eso no depende de los
datos y va \emph{fuera} de \codigo{aes()}:
\codigo{geom\_col(fill = "\#2a78d6")}.
\end{concepto}

\subsection{Los datos del módulo}

Trabajaremos sobre todo con \archivo{transacciones.csv} (1,000 tickets de
2023), unida al catálogo de productos (para tener la categoría) y a las
tiendas (para tener el nombre). Es el mismo \emph{join} que hiciste en el
Módulo~\ref{cap:manipulacion}.

\begin{rcode}
transacciones <- read_csv("datasets/transacciones.csv",
                          show_col_types = FALSE)
productos <- read_csv("datasets/productos.csv", show_col_types = FALSE)
tiendas   <- read_csv("datasets/tiendas.csv", show_col_types = FALSE)

ventas_detalle <- transacciones %>%
  left_join(select(productos, producto_id, categoria),
            by = "producto_id") %>%
  left_join(select(tiendas, tienda_id, nombre_tienda),
            by = "tienda_id") %>%
  # "Sucursal Outlet" -> "Outlet": etiquetas más cortas para las gráficas
  mutate(tienda = str_remove(nombre_tienda, "Sucursal "))

ventas_detalle %>%
  select(fecha, categoria, tienda, metodo_pago, total_transaccion) %>%
  head(4)
\end{rcode}
\begin{rsalida}
PENDIENTE
\end{rsalida}

Para la primera gráfica necesitamos los ingresos de cada mes.
\codigo{floor\_date(fecha, unit = "month")} (de \paquete{lubridate}) lleva
cada fecha al primer día de su mes: así obtenemos una columna de tipo
\emph{fecha}, no de texto, lo que será importante para el eje $x$.

\begin{rcode}
ventas_mes <- ventas_detalle %>%
  mutate(mes = floor_date(fecha, unit = "month")) %>%
  group_by(mes) %>%
  summarise(ingresos = sum(total_transaccion),
            tickets  = n(),
            .groups  = "drop")
ventas_mes
\end{rcode}
\begin{rsalida}
PENDIENTE
\end{rsalida}

\subsection{Tu primer gráfico, capa por capa}

\textbf{Pregunta de negocio:} ¿cómo evolucionaron los ingresos mes a mes en
2023? Construiremos la respuesta en cuatro pasos y guardaremos cada paso en
un objeto para ver cómo cambia.

\begin{rcode}
# Paso 1: datos + estéticas. Solo el "lienzo" con los ejes
paso1 <- ggplot(data = ventas_mes, mapping = aes(x = mes, y = ingresos))
paso1
\end{rcode}

\codigo{ggplot()} crea el lienzo. Con \codigo{aes(x = mes, y = ingresos)}
le dices qué columna va en cada eje. R ya calcula los rangos (de enero a
diciembre; de unos \$57,000 a \$96,000) pero no dibuja nada, porque aún no
hay geometría.

\begin{rcode}
# Paso 2: + una geometría de línea (color y grosor fijos, fuera de aes)
paso2 <- paso1 + geom_line(color = "#2a78d6", linewidth = 0.8)
paso2
\end{rcode}

Al \emph{sumar} \codigo{geom\_line()} aparece la línea. El color y el grosor
(\codigo{linewidth}) son constantes, por eso van como argumentos de la
geometría y no dentro de \codigo{aes()}.

\begin{rcode}
# Paso 3: + puntos en cada mes y + escala del eje y en dólares
paso3 <- paso2 +
  geom_point(color = "#2a78d6", size = 2) +
  scale_y_continuous(labels = dollar)
paso3
\end{rcode}

Las capas se apilan: \codigo{geom\_point()} dibuja encima de la línea. La
\emph{escala} \codigo{scale\_y\_continuous(labels = dollar)} no cambia los
datos, solo cómo se \emph{muestran} los números del eje: 90000 se vuelve
\$90,000.

\begin{rcode}
# Paso 4: + títulos y + un tema limpio
paso4 <- paso3 +
  labs(title = "Ingresos mensuales 2023", x = NULL, y = "Ingresos") +
  theme_minimal() +
  theme(panel.grid.minor = element_blank())
paso4
\end{rcode}

\codigo{labs()} pone los textos (\codigo{x = NULL} quita el título del eje
$x$, porque ``mes'' es obvio). \codigo{theme\_minimal()} quita el fondo gris
y \codigo{theme(panel.grid.minor = element\_blank())} elimina las líneas
secundarias de la cuadrícula, que solo agregan ruido. La
figura~\ref{fig:m03-pasos} muestra los cuatro pasos.

\begin{figure}[H]
\centering
\includegraphics[width=0.48\textwidth]{m03-paso-1}\hfill
\includegraphics[width=0.48\textwidth]{m03-paso-2}\\[4pt]
\includegraphics[width=0.48\textwidth]{m03-paso-3}\hfill
\includegraphics[width=0.48\textwidth]{m03-paso-4}
\caption{La misma gráfica en cuatro pasos: lienzo (arriba a la izquierda),
línea, puntos con escala en dólares y, finalmente, títulos y tema.}
\label{fig:m03-pasos}
\end{figure}

\textbf{Cómo leerla.} Ya con el paso 4 se ve la historia del año: dos caídas
fuertes (febrero y, sobre todo, abril), una meseta alta de mayo a julio, un
bache en septiembre--octubre y un cierre muy fuerte en diciembre.

\begin{advertencia}
El \codigo{+} va \textbf{al final} de cada línea, nunca al principio de la
siguiente. Si escribes el \codigo{+} al inicio, R da por terminada la
instrucción en la línea anterior y luego no sabe qué hacer con el
\codigo{+} suelto. Tampoco uses \codigo{\%>\%} entre capas de
\paquete{ggplot2}: el pipe sirve para \emph{transformar datos}; el \codigo{+}
para \emph{añadir capas}. Verás ambos errores en la
sección~\ref{sec:m3-errores}.
\end{advertencia}

\begin{hazlotu}
Repite los cuatro pasos pero con \codigo{y = tickets} (número de tickets por
mes) y el título ``Tickets por mes 2023''. Pista: en el paso 3 cambia
\codigo{dollar} por \codigo{comma}. ¿Coinciden los meses bajos con los de
ingresos?
\end{hazlotu}

% ----------------------------------------------------------------------------
\section{Elegir la gráfica correcta}
\label{sec:m3-elegir}

Antes de escribir código pregúntate: \emph{¿qué quiero que vea quien mire la
gráfica?} La respuesta casi siempre cae en una de seis familias
(tabla~\ref{tab:m3-elegir}).

\begin{table}[H]
\centering
\small
\caption{¿Qué quieres mostrar? Gráfica recomendada y qué evitar.}
\label{tab:m3-elegir}
\begin{tabularx}{\textwidth}{>{\raggedright\arraybackslash}p{2.6cm}
  >{\raggedright\arraybackslash}X >{\raggedright\arraybackslash}X}
\toprule
\textbf{Quiero mostrar\ldots} & \textbf{Gráfica recomendada} &
\textbf{Evita} \\
\midrule
Comparación entre categorías (tiendas, productos) &
Barras ordenadas; horizontales si los nombres son largos
(\codigo{geom\_col()}) &
Barras de colores sin significado; pasteles; ejes que no empiezan en 0 \\
Evolución en el tiempo (meses, semanas) &
Línea (\codigo{geom\_line()}); barras si son pocos periodos &
Más de 4--5 líneas de colores; doble eje Y \\
Distribución (cómo se reparten los tickets) &
Histograma, densidad, boxplot (\codigo{geom\_histogram()},
\codigo{geom\_boxplot()}) &
Reportar solo el promedio \\
Relación entre dos medidas (precio y cantidad) &
Dispersión (\codigo{geom\_point()}) con tendencia opcional &
Unir los puntos con líneas; más de 3 colores de grupo \\
Composición (partes de un total) &
Barras 100\% o barras de participación; pastel solo con 2--3 porciones &
Pastel de 10 rebanadas; dona 3D \\
Un solo número (ingresos del año) &
Tarjeta KPI: el número grande con su contexto &
Una gráfica con una sola barra \\
\bottomrule
\end{tabularx}
\end{table}

\begin{concepto}{Tarjeta KPI}
Un \termino{KPI} (\emph{Key Performance Indicator}, indicador clave de
desempeño) es un número que resume el desempeño: ingresos, tickets, ticket
promedio. Cuando el mensaje es \emph{un} número, la mejor visualización es
una \termino{tarjeta KPI}: el número grande, una etiqueta y, si acaso, la
comparación contra la meta o el año anterior. Una gráfica con una sola barra
ocupa más espacio y comunica menos.
\end{concepto}

\begin{hazlotu}
Para cada pregunta, decide la familia de la tabla~\ref{tab:m3-elegir}:
(a) ``¿qué vendedor vendió más?'', (b) ``¿cuánto vendimos en total?'',
(c) ``¿los clientes jóvenes gastan más por ticket?'', (d) ``¿cómo se
reparten los ingresos entre los métodos de pago?''.
\end{hazlotu}

